% !TeX program = xelatex
\documentclass{article}
\usepackage{kotex}
\usepackage[a4paper,left=3cm, right=3cm, top=2cm, bottom=2cm]{geometry}
\usepackage{amsmath}
\usepackage{graphicx}
\usepackage{caption}
\usepackage{setspace}
\usepackage{xcolor}
\usepackage{titlesec}
\usepackage{amssymb}
\usepackage{tcolorbox}
\usepackage{wrapfig}
\usepackage{amsthm}
\usepackage{multirow} 
\usepackage{array}
\usepackage{tcolorbox}
\usepackage{pgfplots}
\pgfplotsset{compat=1.18} % For proofs

\graphicspath{{graph/}}
\title{7.5 접평면과 법선}
\date{}
\author{}
\setstretch{1.3}

% \subsection* 형식 지정 (번호 없음)
\titleformat{name=\section, numberless}
  {\normalfont\large\bfseries\color{blue}}
  {}
  {0pt}
  {}
\geometry{a4paper, margin=1in}

% 증명 환경 스타일
\newtheoremstyle{mystyle}% name
  {}% Space above
  {}% Space below
  {\itshape}% Body font
  {}% Indent amount
  {\bfseries}% Theorem head font
  {.}% Punctuation after theorem head
  {.5em}% Space after theorem head
  {}% Theorem head spec (can be left empty, meaning `normal')
\theoremstyle{mystyle}

\begin{document}
\maketitle

\begin{tcolorbox}[colback=blue!5, colframe=blue!60!black, title=7.5 접평면과 법선 -- 한눈에 보기, boxrule=0.5mm, arc=3mm]
곡면 위의 한 점 \(P_0(x_0,y_0,z_0)\)에서,\\ 이차방정식 \(F(x,y,z)=0\)의 \(x^2\to x_0x,\ 2xy\to x_0y+xy_0,\ 2x\to x+x_0\) 식으로 \textbf{절반씩 바꿔치기}(극화, polarization)하면 그 점에서의 접평면 방정식을 바로 얻는다.
\end{tcolorbox}

\begin{tcolorbox}[
    colback=teal!4,
    colframe=teal!65!black,
    title=정의 \& 핵심 규칙,
    fonttitle=\bfseries,
    boxrule=0.5mm,
    arc=3mm
    ]
    \textbf{\textcolor{teal!65!black}{정의.}}\quad
    이차곡면 \(F(x,y,z)=0\) 위의 점 \(P_0(x_0,y_0,z_0)\)를 지나는 모든 접선이 놓이는 평면을 \(P_0\)에서의 \textbf{접평면}(tangent plane)이라 하고, \(P_0\)를 \textbf{접점}(point of tangency)이라 한다. \(P_0\)를 지나고 접평면에 \textbf{수직}인 직선을 \(P_0\)에서의 \textbf{법선}(normal line)이라 한다.

    \noindent\textcolor{teal!50!black}{\rule{\linewidth}{0.4pt}}

    \textbf{\textcolor{teal!65!black}{핵심 규칙(극화).}}\quad
    \(F(x,y,z)=ax^2+by^2+cz^2+2fyz+2gzx+2hxy+2lx+2my+2nz+d\)에서 \\
    $
    x^2\to x_0x,\quad y^2\to y_0y,\quad z^2\to z_0z,\\
    \quad yz\to\tfrac{y_0z+yz_0}{2},\quad zx\to\tfrac{z_0x+zx_0}{2},\quad xy\to\tfrac{x_0y+xy_0}{2},\\
    \quad x\to\tfrac{x+x_0}{2},\ \dots
    $
    로 바꾸면 \(P_0\)에서의 접평면 방정식이 바로 나온다.
\end{tcolorbox}

\textbf{유도.} 점 \(P_0(x_0,y_0,z_0)\)를 지나고 방향코사인이 \(\lambda,\mu,\nu\)인 직선과 곡면의 교점은 (7.3절과 같이)\\
$
(a\lambda^2+b\mu^2+c\nu^2+2f\mu\nu+2g\nu\lambda+2h\lambda\mu)t^2\\
+2\{(ax_0+hy_0+gz_0+l)\lambda+(hx_0+by_0+fz_0+m)\mu+(gx_0+fy_0+cz_0+n)\nu\}t\\
+F(x_0,y_0,z_0)=0
$ \\
을 만족하는 \(t\)로 결정된다. \(P_0\)가 곡면 위에 있으면 \(F(x_0,y_0,z_0)=0\)이므로 \(t=0\)은 항상 근이다. 이 직선이 \(P_0\)에서 \textbf{접한다}는 것은 나머지 한 근도 \(0\)이 된다는 뜻이므로
\[
(ax_0+hy_0+gz_0+l)\lambda+(hx_0+by_0+fz_0+m)\mu+(gx_0+fy_0+cz_0+n)\nu=0
\]
이 조건에서 \(\lambda t=x-x_0,\ \mu t=y-y_0,\ \nu t=z-z_0\)를 이용해 \(\lambda,\mu,\nu\)(따라서 \(t\)도)를 소거하면
\[
(ax_0+hy_0+gz_0+l)(x-x_0)+(hx_0+by_0+fz_0+m)(y-y_0)+(gx_0+fy_0+cz_0+n)(z-z_0)=0
\]
을 얻고, \(F(x_0,y_0,z_0)=0\)을 이용해 정리하면 \textbf{접평면의 방정식}
\[
ax_0x+by_0y+cz_0z+f(y_0z+yz_0)+g(z_0x+zx_0)+h(x_0y+xy_0)+l(x+x_0)+m(y+y_0)+n(z+z_0)+d=0
\]
을 얻는다. 이 평면에 수직이고 \(P_0\)를 지나는 \textbf{법선의 방정식}은
\[
\frac{x-x_0}{ax_0+hy_0+gz_0+l}=\frac{y-y_0}{hx_0+by_0+fz_0+m}=\frac{z-z_0}{gx_0+fy_0+cz_0+n}
\]

% ============ 이미지 삽입 위치: 법선 방정식 직후 -- 접평면과 법선 개념도 ============
% 타원면 위의 한 점 P0에서 접평면(주황 사각형)과 그에 수직인 법선(빨간 직선)을 함께 표시.
\begin{figure}[h]
\centering
\begin{tikzpicture}
\begin{axis}[hide axis, view={115}{20}, width=6cm, height=6cm]
\addplot3[
    surf, shader=faceted, faceted color=teal!55, fill=teal!12, opacity=0.45,
    domain=0:180, domain y=0:360, samples=14, samples y=18, z buffer=sort,
]
({2*sin(x)*cos(y)}, {1.1*sin(x)*sin(y)}, {1.4*cos(x)});
\draw[fill=orange!35, opacity=0.6, draw=orange!70!black, thick]
  (axis cs:0.533,1.478,0.210) -- (axis cs:2.120,0.784,-0.280) --
  (axis cs:2.120,-0.254,1.190) -- (axis cs:0.533,0.441,1.680) -- cycle;
\addplot3[thick, red] coordinates {(1.138,0.324,0.497) (1.752,1.261,1.157)};
\addplot3[only marks, mark=*, mark size=1.3pt, black] coordinates {(1.327,0.612,0.700)};
\node at (axis cs:1.327,0.612,0.700)[anchor=north east]{\footnotesize $P_0$};
\node at (axis cs:1.752,1.261,1.157)[anchor=south west, red]{\footnotesize 법선};
\node at (axis cs:2.120,0.784,-0.280)[anchor=north, orange!70!black]{\footnotesize 접평면};
\end{axis}
\end{tikzpicture}
\caption{타원면 위의 점 $P_0$에서의 접평면(주황)과 법선(빨강) -- 법선은 접평면에 수직}
\end{figure}

\begin{tcolorbox}[colback=gray!5, colframe=gray!55!black, title=보기 1 \& 보기 2, boxrule=0.4mm, arc=2mm]
\textbf{보기 1.} 타원면 \(ax^2+by^2+cz^2=1\) 위의 점 \((x_1,y_1,z_1)\)에서
\[
\text{접평면: } ax_1x+by_1y+cz_1z=1,\qquad
\text{법선: } \frac{x-x_1}{ax_1}=\frac{y-y_1}{by_1}=\frac{z-z_1}{cz_1}
\]

\textbf{보기 2.} 무심이차곡면 \(ax^2+by^2=2cz\) 위의 점 \((x_1,y_1,z_1)\)에서
\[
\text{접평면: } ax_1x+by_1y=c(z+z_1),\qquad
\text{법선: } \frac{x-x_1}{ax_1}=\frac{y-y_1}{by_1}=\frac{z-z_1}{-c}
\]
\end{tcolorbox}

\subsection*{EXAMPLE 1}
방향코사인이 \(\lambda,\mu,\nu\)이고 타원면 \(\dfrac{x^2}{a^2}+\dfrac{y^2}{b^2}+\dfrac{z^2}{c^2}=1\)에 접하는 접평면의 방정식은
\[
\lambda x+\mu y+\nu z=\sqrt{a^2\lambda^2+b^2\mu^2+c^2\nu^2}
\]
임을 증명하여라.\\
\textbf{SOLUTION:} 이 타원면 위의 점 \(P_0(x_0,y_0,z_0)\)에서의 접평면은 \(\dfrac{x_0x}{a^2}+\dfrac{y_0y}{b^2}+\dfrac{z_0z}{c^2}=1\)이고, 방향코사인이 \(\lambda,\mu,\nu\)인 접평면을 \(\lambda x+\mu y+\nu z=p\)라 하면 두 식이 일치하려면
\[
\frac{x_0/a^2}{\lambda}=\frac{y_0/b^2}{\mu}=\frac{z_0/c^2}{\nu}=\frac{1}{p}
\quad\Longrightarrow\quad
x_0=\frac{a^2\lambda}{p},\quad y_0=\frac{b^2\mu}{p},\quad z_0=\frac{c^2\nu}{p}
\]
이를 \(\dfrac{x_0^2}{a^2}+\dfrac{y_0^2}{b^2}+\dfrac{z_0^2}{c^2}=1\)에 대입하면
\[
\frac{a^2\lambda^2+b^2\mu^2+c^2\nu^2}{p^2}=1
\quad\Longrightarrow\quad
p=\pm\sqrt{a^2\lambda^2+b^2\mu^2+c^2\nu^2}
\]
그러므로 구하는 접평면의 방정식은 \(\lambda x+\mu y+\nu z=\sqrt{a^2\lambda^2+b^2\mu^2+c^2\nu^2}\)이다.

같은 방법으로 다음도 얻는다.
\begin{center}
\renewcommand{\arraystretch}{1.4}
\begin{tabular}{|l|l|}
\hline
곡면 & 방향코사인 $\lambda,\mu,\nu$인 접평면 \\
\hline
일엽쌍곡면 \(\dfrac{x^2}{a^2}+\dfrac{y^2}{b^2}-\dfrac{z^2}{c^2}=1\) & \(\lambda x+\mu y+\nu z=\sqrt{a^2\lambda^2+b^2\mu^2-c^2\nu^2}\) \\
\hline
이엽쌍곡면 \(\dfrac{z^2}{c^2}-\dfrac{x^2}{a^2}-\dfrac{y^2}{b^2}=1\) & \(\lambda x+\mu y+\nu z=\sqrt{c^2\nu^2-a^2\lambda^2-b^2\mu^2}\) \\
\hline
\end{tabular}
\end{center}

\section*{연습문제 7.5}
\begin{enumerate}
    \item 이차곡면 \(ax^2+by^2+cz^2=0\)에 접하고 방향코사인이 \(\lambda,\mu,\nu\)인 접평면의 방정식과, 곡면 위의 점 \(P(x_0,y_0,z_0)\)에서의 법선의 방정식을 구하여라.
    \item 이차곡면 \(ax^2+by^2=2z\)에 평면 \(lx+my+nz=p\)가 접할 조건을 구하여라.
    \item 타원면 \(\dfrac{x^2}{a^2}+\dfrac{y^2}{b^2}+\dfrac{z^2}{c^2}=1\) 위의 한 점 \(P\)에서의 법선이 \(yz,\ zx,\ xy\)평면과 만나는 점을 각각 \(A,B,C\)라 하면
    \[
    \overline{PA}:\overline{PB}:\overline{PC}=a^2:b^2:c^2
    \]
    임을 증명하여라.
\end{enumerate}
\end{document}